\section{Overall Description}
\label{sec:overall-description}

\subsection{Product Perspective}
Campus ERP is a new, self-contained system composed of four independently
deployable services behind one API gateway (\cref{fig:d5-deployment}). It does
not integrate with any pre-existing institutional system; it is designed so
that each domain service could, in principle, be developed, tested, deployed
and scaled by a separate team without touching the others' code or database.

\subsection{Product Functions}
\begin{itemize}[leftmargin=1.4em]
  \item \textbf{Identity \& Access}: multi-tenant sign-in, role-based
    permissions, session and refresh-token lifecycle, tenant and user
    administration.
  \item \textbf{Academic}: programme/course catalogue with prerequisites,
    term registration, gradebook, attendance, conflict-checked examination
    scheduling, result publication, grade appeals, at-risk flagging,
    transcript/attendance PDF export.
  \item \textbf{Marketing \& Finance}: fee plans, tuition invoicing (driven by
    Academic registration events), simulated mobile-money payments, digital
    receipts, expense approval workflow, marketing campaign/lead tracking with
    ROI, automated monthly financial reporting, general ledger.
  \item \textbf{Administration \& HR}: employee records, recruitment pipeline,
    statutory payroll (CNPS/IRPP), QR-based attendance, leave management,
    performance reviews, asset and inventory tracking, notifications.
\end{itemize}

\subsection{User Classes and Characteristics}
\begin{longtable}{L{2.6cm}L{11.5cm}}
  \toprule
  \rowcolor{ceTeal}\thd{Role} & \thd{Characteristics and typical tasks} \\
  \midrule
  \endhead
  Super Admin & Platform-level operator; provisions institutions (tenants) and
    their first administrator. Does not gain implicit access to any tenant's
    business data. \\
  Admin & Institution administrator; manages the academic catalogue, users and
    roles within their own tenant. \\
  Instructor & Manages courses they are assigned to: gradebook, attendance,
    exam scheduling, appeal review. Cannot see or modify other instructors'
    courses. \\
  Finance Officer & Manages fee plans, payments, expenses, campaigns and
    reports for their tenant. \\
  HR Officer & Manages employees, recruitment, payroll, leave, performance and
    assets for their tenant. \\
  Employee & Self-service: QR check-in/out, leave requests, payslips,
    performance history -- limited to their own records. \\
  Student & Self-service: course results, exam timetable, grade appeals,
    tuition payment, transcript download -- limited to their own records. \\
  \bottomrule
\end{longtable}

\subsection{Operating Environment}
The system runs as Docker containers orchestrated by \texttt{docker compose}
on a single host for this course deliverable (a streaming standby and a
horizontal-scaling story are documented conceptually in
\cref{sec:nfr-testing-devops}). The frontend is a modern evergreen browser;
the API is consumed over HTTPS. All displayed dates and times use the
Africa/Douala timezone (WAT, UTC+1); all money values are stored as integer
FCFA and displayed accordingly.

\subsection{Design and Implementation Constraints}
The brief's mandatory architectural constraints, and how Campus ERP satisfies
each, are summarised in \cref{tab:constraints}.

\begin{table}[H]
  \centering
  \caption{Mandatory architectural constraints and how they are satisfied}
  \label{tab:constraints}
  \begin{tabularx}{\linewidth}{L{5.1cm}X}
    \toprule
    \rowcolor{ceTeal}\thd{Constraint} & \thd{Implementation} \\
    \midrule
    $\geq$3 independently deployable microservices, database-per-service &
      Four services (Identity, Academic, Finance, HR), each with its own
      PostgreSQL database and least-privilege runtime role. \\
    API Gateway: routing, auth delegation, rate limiting &
      NGINX terminates TLS, routes \texttt{/api/v1/*} by prefix, delegates
      session validation via \texttt{auth\_request}, and enforces per-IP rate
      limits (general traffic and a stricter limit on login/refresh). \\
    Asynchronous messaging for $\geq$1 cross-module workflow &
      RabbitMQ topic exchange \texttt{campus.events}; Academic publishes
      \texttt{academic.enrollment.confirmed.v1} via a transactional outbox,
      Finance consumes it idempotently to create the tuition invoice
      (\cref{fig:d3a-sequence-enrollment-invoice}). \\
    Containers + \texttt{docker-compose}, one-command startup &
      A Dockerfile per service (and per worker), a gateway image bundling the
      built frontend, and one \texttt{docker-compose.yml} bringing up the
      full stack including PostgreSQL, RabbitMQ, observability and backup
      services. \\
    CI/CD pipeline running tests and building images on every push to main &
      GitHub Actions workflow: lint, per-service test + coverage gate, frontend
      build, then multi-service image build on \texttt{main}. \\
    \bottomrule
  \end{tabularx}
\end{table}

\subsection{Assumptions and Dependencies}
\begin{itemize}[leftmargin=1.4em]
  \item Mobile-money payment is demonstrated with a disclosed simulator that
    mimics the MTN MoMo / Orange Money request-to-pay and signed-callback
    flow; no real sandbox credentials were available, and the brief
    explicitly permits a mocked equivalent.
  \item Statutory payroll rates (CNPS, IRPP/PAYE, CAC) are implemented as an
    editable, source-cited rule set that must be reviewed and approved by an
    HR officer before payroll can run -- see the governance note in
    \cref{sec:hr-module} and sources in Appendix~\ref{app:sources}.
  \item The local deployment target is a single Docker host; a streaming
    standby profile is included for the backup/recovery story, but automatic
    failover is out of scope.
  \item All demonstration data (institutions, staff, students) is synthetic.
\end{itemize}

\AppFigure{login-page}{0.92}{Sign-in screen, showing the light mist/teal/coral design and the demonstration-account picker}
